\documentclass[a4paper,11pt]{article}
\usepackage{graphicx} % Required for inserting images
\usepackage{float}
\usepackage{amsmath}
\usepackage{amssymb}
\usepackage{textalpha}
\usepackage{enumitem}
\usepackage{changepage} % adjustwidth
\usepackage[utf8]{inputenc}  
\usepackage{titlesec}
\usepackage[LGR,T1]{fontenc} 
\usepackage[top=3cm,bottom=3cm,left=2cm,right=2cm,marginparwidth=1.75cm]{geometry}


\usepackage[greek,english]{babel}  


\pagestyle{empty}

\begin{document}
\renewcommand{\figurename}{Σχήμα}
\renewcommand{\tablename}{Πίνακας}

\textbf{Θεωρητικό Υπόβαθρο}\\

\textbf{1. Mel-frequency Cepstral Coefficients (MFCCs)}\\

\textbf{2. Γλωσσικά Μοντέλα (Language Models)}\\

\textgreek{Ως γλωσσικό μοντέλο ορίζουμε ένα μοντέλο μηχανικής μάθησης, το οποίο εκτιμά την πιθανότητα εμφάνισης μελλοντικών λέξεων σε ένα κείμενο. Δηλαδή, πρόκειται για ένα μοντέλο που σύμφωνα με κάποιο κριτήριο υπολογίζει την πιθανότητα μία λέξη ή μία πρόταση να ακολουθεί του τωρινού κειμένου. Για τον σκοπό αυτό, απαιτείται συνήθως το σύνολο των λέξεων, το λεξικό, να είναι γνωστό.}

\textgreek{Τα γλωσσικά μοντέλα είναι ένα βασικό εργαλείο με το οποίο δύναται να εκτιμηθεί η συνοχή και η συνέπεια ενός κειμένου, ώστε να εντοπιστούν και να διορθωθούν τυχόν λάθη. Σε ένα κείμενο, μπορεί να υπάρχουν διάφοροι τύποι σφαλμάτων διαφορετικής δυσκολίας εντοπισμού. Λάθη γραμματικής, ορθογραφικά και τυπογραφικά, εντοπίζονται με μία αναζήτηση στο λεξικό και μπορούν να διορθωθούν από σχετικά απλά γλωσσικά μοντέλα. Ωστόσο, ασάφειες και αντιφάσεις είναι πιο δύσκολο να αντιμετωπιστούν και συχνά απαιτούν προηγμένα γλωσσικά μοντέλα για να ανιχνευθούν και να διορθωθούν.}

\textgreek{Συνεπώς, τα γλωσσικά μοντέλα μπορούν να περιορίσουν τα σφάλματα του κειμένου εισόδου και βρίσκουν εφαρμογή σε συστήματα επεξεργασίας φυσικής γλώσσας.}

\textgreek{Υπάρχουν διάφορες κατηγορίες γλωσσικών μοντέλων, ορισμένα πιο απλά ως προς την υλοποίησή τους, όπως τα }N-gram language models \textgreek{ και άλλα περισσότερο περίπλοκα, όπως τα }Large language models. \textgreek{Στην παρούσα εργασία όμως θα ασχοληθούμε με τα} N-gram language models, \textgreek{τα οποία ορίζονται ως εξής.}

\textgreek{Έστω }$w_{i+1}, w_{i+2}, w_{i+3}, ..., w_{i+N-1}$ \textgreek{οι τελευταίες }$N-1$ \textgreek{ λέξεις του κειμένου. Τότε, η πιθανότητα η επόμενη λέξη του κειμένου να είναι η λέξη }$x$, \textgreek{είναι ίση με:}

\[\Pr[w_{i+N}=x]=\Pr[x \mid w_{i+1}w_{i+2}...w_{i+N-1}]\]

\textgreek{Τα }N-gram language models\textgreek{ είναι χρήσιμα σε περιπτώσεις όπου τα }Large language models \textgreek{είναι δύσκολα υπολογιστικά, ωστόσο δεν χειρίζονται αποτελεσματικά τις συσχετίσεις μεταξύ λέξεων που εμφανίζουν απόσταση μεταξύ τους και δεν μοντελοποιούν ικανοποιητικά μη γνωστές ακολουθίες λέξεων. }

\textgreek{Στην άσκηση, θα εργαστούμε με }unigram ($N=1$) \textgreek{και }bigram ($N=2$) \textgreek{μοντέλα.}

\textgreek{Στην περίπτωση των }unigram \textgreek{μοντέλων, οι πιθανότητες μπορούν να υπολογιστούν εύκολα ως } $\Pr[w]=\dfrac{count(w)}{total \ words}$, \textgreek{ενώ για τα }bigram \textgreek{μοντέλα} $\Pr[w_{i} \mid w_{i-1}]=\dfrac{\Pr[w_{i-1}w_{i}]}{\Pr[w_{i-1}]}=\dfrac{count(w_{i-1}w_{i})}{count(w_{i-1})}$. \textgreek{Βέβαια, για την αποφυγή }underflow\textgreek{ λόγω συνεχόμενων πολλαπλασιασμών με μικρούς αριθμούς, οι πιθανότητες υπολογίζονται και αποθηκεύονται σε λογαριθμική μορφή.}

\textgreek{Ωστόσο, τα }N-gram language models \textgreek{παρουσιάζουν καλές επιδόσεις μόνο εάν το }test set \textgreek{είναι πανομοιότυπο με το }train set. \textgreek{Ακόμα βέβαια και σε αυτές τις περιπτώσεις, δύναται να υπάρξει στο }test set \textgreek{κάποια ακολουθία λέξεων που δεν παρουσιάζεται στο }train set, \textgreek{με αποτέλεσμα το μοντέλο μας να έχει αναθέσει λανθασμένα μηδενική πιθανότητα στη συγκεκριμένη ακολουθία. Αυτός είναι και ο βασικός λόγος για τον οποίο τα }N-gram language models \textgreek{δεν παρουσιάζουν σθεναρότητα και δεν είναι κατάλληλα για γενίκευση.}

\textgreek{Προκειμένου να αποφύγουμε τις μηδενικές πιθανότητες, αναθέτουμε με βάση κάποιο κριτήριο (π.χ. }add-one, interpolation, backoff) \textgreek{μικρή πιθανότητα στις ακολουθίες που δεν παρουσιάζονται στο }training set.

\textgreek{Για την αξιολόγηση των γλωσσικών μοντέλων θα πρέπει να εκπαιδεύσουμε τα μοντέλα σε ένα }train \textgreek{κείμενο να αξιολογήσουμε την επίδοσή τους σε ένα }test \textgreek{κείμενο σύμφωνα με κάποια μετρική. Μία τέτοια μετρική είναι η }perplexity (PP).

\[
PP(w) = 2^{H(w)}, \ H(w) =− \dfrac{1}{M} log_2P(w)
\]

\textgreek{όπου Μ το πλήθος των λέξεων. Μάλιστα, ισχύει ότι}

\[
PP(w) = 2^{H(w)} = P(w_1, w_2, \ldots, w_M)^{-\frac{1}{M}}, \quad \text{καθώς } M \to \infty
\]

\textgreek{Σύμφωνα με τη συγκεκριμένη μετρική, όταν οι προβλέψεις των λέξεων γίνονται πιο ακριβείς, δηλαδή όταν αυξάνονται οι δεσμευμένες πιθανότητες, η τιμή }perplexity \textgreek{ελαττώνεται. Συνεπώς, ένα μοντέλο αξιολογείται ως καλύτερο από ένα άλλο αν εμφανίζει μικρότερη τιμή }perplexity. 

\textgreek{Στην περίπτωση του }unigram \textgreek{μοντέλου} $
H(w) = -\frac{1}{M} \sum_{n=1}^{M} \log_2 P(w_n)$, \textgreek{ενώ στην περίπτωση του }bigram \textgreek{ μοντέλου} $H(w) = -\frac{1}{M} \sum_{n=1}^{M} \log_2 P(w_n \mid w_{n-1})$.

\textgreek{Η }perplexity \textgreek{είναι μία }intrinsic evaluation, \textgreek{χρήσιμη για την εξαγωγή συμπερασμάτων για την απόδοση των γλωσσικών μοντέλων, ωστόσο δεν είναι ασφαλής για την εκτίμηση της επίδοσης τους σε πραγματικές εφαρμογές. }


\end{document}
